\label{kap:absolutni}
\mathquote{Mathematics is not about numbers, equations, computations, or algorithms: it is about understanding.}{William Paul Thurston}

V této kapitole se podíváme na dva typy rovnic, které vypadají nevinně, ale umí potrápit. Rovnice s absolutní hodnotou, u kterých musíme rozebírat případy, a iracionální rovnice, tedy rovnice s neznámou pod odmocninou, u kterých se poprvé setkáme s tím, že správně provedený výpočet vyprodukuje špatné řešení. Obojí spojuje jedna věc: musíme hodně dávat pozor na to, co přesně děláme.

\section{Absolutní hodnota}

\begin{definitionbox}
  \textbf{Absolutní hodnota} reálného čísla $a$ je číslo
  \begin{equation*}
    |a| = \begin{cases}
      a & \text{pro } a\geq 0,\\
      -a & \text{pro } a < 0.
    \end{cases}
  \end{equation*}
  \textbf{Geometrický význam}: $|a|$ je vzdálenost čísla $a$ od nuly na číselné ose. Obecněji $|a - b|$ je \textbf{vzdálenost} čísel $a$ a $b$.
\end{definitionbox}

Definice na první pohled děsí tím, že se v ní objevuje $-a$. Absolutní hodnota přece nemůže být záporná! A taky není. Pokud je $a < 0$, je $-a$ kladné číslo. Například $|-5| = -(-5) = 5$. Absolutní hodnota z čísla prostě udělá kladné číslo, bez ohledu na to, čím si prošlo.

\begin{center}
  \begin{tikzpicture}
    \draw[->, thick] (-4.5,0) -- (4.5,0) node[right] {$x$};
    \foreach \x in {-4,...,4} \draw (\x,0.1) -- (\x,-0.1) node[below] {\small $\x$};
    \draw[colorofwisdom, line width=2pt] (-3,0.4) -- (0,0.4);
    \node[above, colorofwisdom] at (-1.5,0.4) {$|-3| = 3$};
    \draw[line width=2pt] (1,-0.8) -- (4,-0.8);
    \node[below] at (2.5,-0.8) {$|4 - 1| = |1 - 4| = 3$};
  \end{tikzpicture}
\end{center}

Jak vypadá graf funkce $y = |x|$? Pro $x\geq 0$ je to přímka $y = x$, pro $x < 0$ přímka $y = -x$. Jinak řečeno: vezmeme přímku $y = x$ a její část pod osou $x$ \textit{překlopíme} nahoru. Vznikne tvar písmene V se špičkou v počátku. Posouvání pak funguje jako u každé jiné funkce. Graf $y = |x - 2|$ má špičku v bodě $2$, tedy tam, kde je výraz uvnitř absolutní hodnoty nulový, a graf $y = |x| - 1$ je celé V posunuté o $1$ dolů:
\begin{center}
  \begin{tikzpicture}[scale=0.55]
    \begin{scope}
      \draw[->] (-2.8,0) -- (3.3,0) node[right] {\small $x$};
      \draw[->] (0,-2) -- (0,3.4) node[above] {\small $y$};
      \draw[dashed, gray, thick] (-1.8,-1.8) -- (0,0);
      \draw[->, gray] (-1.3,-1.1) to[bend left=40] (-1.3,1.1);
      \draw[very thick, colorofwisdom] (-2.6,2.6) -- (0,0) -- (2.9,2.9);
      \node[font=\small] at (0.3,-2.7) {$y = |x|$};
    \end{scope}
    \begin{scope}[xshift=7.3cm]
      \draw[->] (-1.4,0) -- (5,0) node[right] {\small $x$};
      \draw[->] (0,-2) -- (0,3.4) node[above] {\small $y$};
      \draw[very thick, colorofwisdom] (-1,3) -- (2,0) -- (4.8,2.8);
      \filldraw (2,0) circle (3pt) node[below, font=\small] {$2$};
      \node[font=\small] at (1.8,-2.7) {$y = |x - 2|$};
    \end{scope}
    \begin{scope}[xshift=15.5cm]
      \draw[->] (-2.8,0) -- (3.3,0) node[right] {\small $x$};
      \draw[->] (0,-2) -- (0,3.4) node[above] {\small $y$};
      \draw[dashed, gray] (-2.6,2.6) -- (0,0) -- (2.9,2.9);
      \draw[very thick, colorofwisdom] (-2.6,1.6) -- (0,-1) -- (2.9,1.9);
      \filldraw (-1,0) circle (3pt) node[below left, font=\small] {$-1$};
      \filldraw (1,0) circle (3pt) node[below right, font=\small] {$1$};
      \node[font=\small] at (0.3,-2.7) {$y = |x| - 1$};
    \end{scope}
  \end{tikzpicture}
\end{center}
Ostrá špička je typickým poznávacím znamením absolutní hodnoty. V kapitole o derivacích uvidíme, že právě v ní nemá graf tečnu.

\begin{definitionbox}
  \textbf{Vlastnosti absolutní hodnoty}: Pro všechna $a, b\in\mathbb{R}$ platí
  \begin{enumerate}
    \item $|a|\geq 0$ a $|a| = 0\Leftrightarrow a = 0$,
    \item $|-a| = |a|$,
    \item $|a\cdot b| = |a|\cdot|b|$ a $\left|\frac{a}{b}\right| = \frac{|a|}{|b|}$ pro $b\neq 0$,
    \item $|a|^2 = a^2$ a $\sqrt{a^2} = |a|$,
    \item $-|a|\leq a\leq|a|$,
    \item $|a + b|\leq|a| + |b|$ \quad (\textbf{trojúhelníková nerovnost}).
  \end{enumerate}
\end{definitionbox}

Většina vlastností plyne přímo z definice rozebráním případů $a\geq 0$ a $a < 0$. Zkuste si to. Nejzajímavější je ta poslední, a tak si ji dokážeme.

\begin{proof}[(trojúhelníkové nerovnosti)]
  Obě strany jsou nezáporné, takže stačí dokázat nerovnost pro jejich druhé mocniny.\footnote{Pro nezáporná čísla $u, v$ platí $u\leq v\Leftrightarrow u^2\leq v^2$. Pro záporná čísla by to neplatilo, viz $-3\leq 1$, ale $9 > 1$.} Podle vlastnosti 4 je
  \begin{equation*}
    |a + b|^2 = {(a + b)}^2 = a^2 + 2ab + b^2,\qquad {(|a| + |b|)}^2 = a^2 + 2|a||b| + b^2 = a^2 + 2|ab| + b^2.
  \end{equation*}
  Podle vlastnosti 5 je $ab\leq|ab|$, a tedy $|a + b|^2\leq{(|a| + |b|)}^2$.
\end{proof}

Proč se jí říká trojúhelníková? Ve vektorech (kapitola o analytické geometrii) bude mít tvar $|\vec{u} + \vec{v}|\leq|\vec{u}| + |\vec{v}|$ a říká, že žádná strana trojúhelníku není delší než součet zbylých dvou. Neboli: přímá cesta je vždy nejkratší. Rovnost nastává, právě když $a$ a $b$ mají stejné znaménko (nebo je jedno z nich nula).

Už na číselné ose je ale vidět, o co jde. Číslo $a + b$ najdeme tak, že z nuly uděláme krok $a$ a pak krok $b$. Celkem jsme ušli $|a| + |b|$. Pokud oba kroky míří stejným směrem, skončíme přesně v této vzdálenosti od nuly. Pokud míří proti sobě, druhý krok část cesty `vrátí' a skončíme blíž:
\begin{center}
  \begin{tikzpicture}[scale=0.85]
    \begin{scope}
      \draw[->] (-1,0) -- (4.5,0);
      \foreach \x in {0,...,4} \draw (\x,0.08) -- (\x,-0.08) node[below, font=\scriptsize] {$\x$};
      \draw[->, very thick, colorofwisdom] (0,0.3) -- (2,0.3) node[midway, above, font=\small] {$a = 2$};
      \draw[->, very thick] (2,0.3) -- (3.5,0.3) node[midway, above, font=\small] {$b = 1{,}5$};
      \filldraw (3.5,0) circle (2.5pt);
      \node[font=\small] at (1.75,-1) {$|a + b| = 3{,}5 = |a| + |b|$};
    \end{scope}
    \begin{scope}[xshift=7.5cm]
      \draw[->] (-1,0) -- (4.5,0);
      \foreach \x in {0,...,4} \draw (\x,0.08) -- (\x,-0.08) node[below, font=\scriptsize] {$\x$};
      \draw[->, very thick, colorofwisdom] (0,0.3) -- (3,0.3) node[midway, above, font=\small] {$a = 3$};
      \draw[->, very thick] (3,1.05) -- (1,1.05) node[midway, above, font=\small] {$b = -2$};
      \draw[dotted] (3,0.3) -- (3,1.05);
      \filldraw (1,0) circle (2.5pt);
      \node[font=\small] at (1.75,-1) {$|a + b| = 1 < 5 = |a| + |b|$};
    \end{scope}
  \end{tikzpicture}
\end{center}

\begin{problem}
  \begin{enumerate}[label=(\alph*)]
    \item Spočítejte $|3 - \pi|$, $|\sqrt{2} - 1|$ a $|1 - \sqrt{3}|$.
    \item Zjednodušte $|x - 2| + |x + 1|$ pro $x\in(-1, 2)$.
    \item Najděte $a, b$ taková, aby $|a + b| < |a| + |b|$.
  \end{enumerate}
\end{problem}

\begin{solution}
  \begin{enumerate}[label=(\alph*)]
    \item $\pi > 3$, takže $|3 - \pi| = \pi - 3$. Dále $\sqrt{2} > 1$, takže $|\sqrt{2} - 1| = \sqrt{2} - 1$, a $\sqrt{3} > 1$, takže $|1 - \sqrt{3}| = \sqrt{3} - 1$.
    \item Pro $x\in(-1, 2)$ je $x - 2 < 0$ a $x + 1 > 0$, takže výraz je $-(x - 2) + (x + 1) = 3$. Pro všechna taková $x$ vychází stejně. To je geometricky jasné: součet vzdáleností bodu mezi $-1$ a $2$ od těchto dvou bodů je přesně vzdálenost mezi nimi.
    \item Stačí čísla s opačným znaménkem, například $a = 1$, $b = -1$: $0 < 2$.
  \end{enumerate}
\end{solution}

\section{Rovnice a nerovnice s absolutní hodnotou}

\subsection{Jednoduché případy}

Pokud je v rovnici jen jedna absolutní hodnota tvaru $|x - a|$, nejrychlejší je myslet geometricky.

\begin{definitionbox}
  Pro $b > 0$ platí:
  \begin{align*}
    |x - a| = b &\quad\Leftrightarrow\quad x = a - b\ \lor\ x = a + b,\\
    |x - a| < b &\quad\Leftrightarrow\quad a - b < x < a + b,\\
    |x - a| > b &\quad\Leftrightarrow\quad x < a - b\ \lor\ x > a + b.
  \end{align*}
  Slovy: hledáme čísla, jejichž vzdálenost od $a$ je rovna, menší nebo větší než $b$. Pro $b < 0$ nemá rovnice $|x - a| = b$ řešení (vzdálenost nemůže být záporná) a pro $b = 0$ má jediné řešení $x = a$.
\end{definitionbox}

Představte si, že do bodu $a$ zapíchneme kružítko s rozevřením $b$. Body, kam kružítko dosáhne přesně, řeší rovnici. Body uvnitř řeší nerovnici `$<$' a body venku nerovnici `$>$':
\begin{center}
  \begin{tikzpicture}
    \draw[->, thick] (-4.5,0) -- (4.5,0) node[right] {$x$};
    \draw[line width=3pt, colorofwisdom] (-2.5,0) -- (2.5,0);
    \draw[line width=3pt, black!25] (-4.3,0) -- (-2.5,0);
    \draw[line width=3pt, black!25] (2.5,0) -- (4.3,0);
    \filldraw (0,0) circle (2pt);
    \filldraw[fill=white, draw=black, thick] (-2.5,0) circle (3pt);
    \filldraw[fill=white, draw=black, thick] (2.5,0) circle (3pt);
    \node[below] at (0,-0.15) {$a$};
    \node[below] at (-2.5,-0.15) {$a - b$};
    \node[below] at (2.5,-0.15) {$a + b$};
    \draw[decorate, decoration={brace, amplitude=5pt}] (-2.5,0.2) -- (0,0.2) node[midway, above=4pt, font=\small] {$b$};
    \draw[decorate, decoration={brace, amplitude=5pt}] (0,0.2) -- (2.5,0.2) node[midway, above=4pt, font=\small] {$b$};
    \node[colorofwisdom, font=\small] at (0,-0.95) {$|x - a| < b$};
    \node[gray, font=\small] at (-3.6,-0.95) {$|x - a| > b$};
    \node[gray, font=\small] at (3.6,-0.95) {$|x - a| > b$};
    \node[font=\small] at (0,1.25) {krajní body: $|x - a| = b$};
  \end{tikzpicture}
\end{center}

\begin{example}
  Vyřešme:
  \begin{enumerate}[label=(\alph*)]
    \item $|x - 3| = 5$. Čísla ve vzdálenosti $5$ od $3$ jsou $-2$ a $8$. $K = \{-2, 8\}$.
    \item $|x + 2|\leq 1$. Protože $x + 2 = x - (-2)$, hledáme čísla ve vzdálenosti nejvýše $1$ od $-2$. $K = [-3, -1]$.
    \item $|2x - 1| > 3$. Buď $2x - 1 > 3$, tedy $x > 2$, nebo $2x - 1 < -3$, tedy $x < -1$. $K = (-\infty, -1)\cup(2, \infty)$.
    \item $|x - 7| = -2$. Absolutní hodnota nemůže být záporná, $K = \emptyset$. Nic nepočítáme.
  \end{enumerate}
\end{example}

\subsection{Metoda nulových bodů}

Co když je v rovnici absolutních hodnot víc, nebo je kombinovaná s dalšími výrazy? Pak použijeme obecnou metodu. Myšlenka je jednoduchá: absolutní hodnota je `problém' jen proto, že nevíme, jestli je výraz uvnitř kladný, nebo záporný. Rozdělíme tedy osu na intervaly, na kterých to víme.

\begin{definitionbox}
  \textbf{Metoda nulových bodů pro absolutní hodnoty}:
  \begin{enumerate}
    \item Najdeme \textbf{nulové body} všech výrazů v absolutních hodnotách.
    \item Nulové body rozdělí reálnou osu na intervaly. Na každém intervalu má každý výraz stálé znaménko, takže absolutní hodnoty můžeme \textbf{odstranit} (výraz buď opíšeme, nebo vezmeme s opačným znaménkem).
    \item Na každém intervalu zvlášť vyřešíme vzniklou rovnici bez absolutních hodnot.
    \item \textbf{Ověříme}, zda nalezené řešení leží v daném intervalu. Pokud ne, zahodíme ho.
    \item Výsledek je sjednocení řešení ze všech intervalů.
  \end{enumerate}
\end{definitionbox}

Krok 4 je důležitý a zapomíná se na něj nejčastěji. Na každém intervalu řešíme jinou rovnici, která platí \textit{jen} na tom intervalu. Řešení mimo interval je řešením špatné rovnice.

\begin{example}\label{ex:abs-dve}
  Vyřešme $|x + 1| + |x - 2| = 5$. Nulové body jsou $-1$ a $2$. Sestavíme tabulku (do které bodu nulové body přiřadíme, je jedno, jen musíme být konzistentní):
  \begin{center}
    \renewcommand{\arraystretch}{1.4}
    \begin{tabular}{|c|c|c|c|}
      \hline
       & $(-\infty, -1)$ & $[-1, 2)$ & $[2, \infty)$ \\
      \hline
      $|x + 1|$ & $-x - 1$ & $x + 1$ & $x + 1$ \\
      $|x - 2|$ & $-x + 2$ & $-x + 2$ & $x - 2$ \\
      \hline
      rovnice & $-2x + 1 = 5$ & $3 = 5$ & $2x - 1 = 5$ \\
      řešení & $x = -2$ & --- & $x = 3$ \\
      leží v intervalu? & ano & --- & ano \\
      \hline
    \end{tabular}
  \end{center}
  $K = \{-2, 3\}$. Geometricky: hledáme body, jejichž součet vzdáleností od $-1$ a od $2$ je $5$. Body mezi $-1$ a $2$ mají tento součet vždy $3$, takže musíme jít ven. A vyjít musíme o $1$ na každou stranu.
\end{example}

Celý příklad si můžeme nakreslit. Graf funkce $y = |x + 1| + |x - 2|$ se skládá ze tří úseček, na každém intervalu z tabulky jedné. Nulové body jsou přesně místa, kde se graf `láme'. Uprostřed je graf vodorovný ve výšce $3$, po stranách stoupá. Řešení rovnice jsou $x$-ové souřadnice průsečíků s vodorovnou přímkou $y = 5$:
\begin{center}
  \begin{tikzpicture}[xscale=0.85, yscale=0.5]
    \draw[->] (-3.8,0) -- (4.8,0) node[right] {\small $x$};
    \draw[->] (0,-0.5) -- (0,7.8) node[above] {\small $y$};
    \foreach \x in {-3,...,4} \draw (\x,0.12) -- (\x,-0.12) node[below, font=\scriptsize] {$\x$};
    \draw[dashed, gray] (-1,0) -- (-1,3);
    \draw[dashed, gray] (2,0) -- (2,3);
    \draw[very thick, colorofwisdom] (-3,7) -- (-1,3) -- (2,3) -- (4,7);
    \draw[thick] (-3.6,5) -- (4.6,5) node[right, font=\small] {$y = 5$};
    \filldraw (-2,5) circle (3pt);
    \filldraw (3,5) circle (3pt);
    \draw[dotted, thick] (-2,5) -- (-2,0);
    \draw[dotted, thick] (3,5) -- (3,0);
    \node[left, font=\small] at (-3.05,7) {$-2x + 1$};
    \node[above, font=\small] at (0.5,3) {$3$};
    \node[right, font=\small] at (4.05,7) {$2x - 1$};
  \end{tikzpicture}
\end{center}
Z obrázku je hned vidět i to, že rovnice $|x + 1| + |x - 2| = c$ nemá pro $c < 3$ žádné řešení.

\begin{example}
  Vyřešme $|x - 1| = 2x + 4$. Nulový bod je $1$.
  \begin{itemize}
    \item Pro $x < 1$: $-x + 1 = 2x + 4$, tedy $x = -1$. To leží v $(-\infty, 1)$, takže je to řešení.
    \item Pro $x\geq 1$: $x - 1 = 2x + 4$, tedy $x = -5$. To \textbf{neleží} v $[1, \infty)$, zahodíme.
  \end{itemize}
  $K = \{-1\}$. Zkouška: $|-2| = 2 = 2\cdot(-1) + 4$. Kdybychom zapomněli na krok 4, tvrdili bychom, že $x = -5$ je řešení, ale $|-6| = 6\neq -6$.
\end{example}

Odkud se `řešení' $x = -5$ vzalo? Na intervalu $[1, \infty)$ jsme řešili rovnici $x - 1 = 2x + 4$. Tím jsme vlastně hledali průsečík přímky $y = 2x + 4$ s \textit{celou} přímkou $y = x - 1$, a ne jen s její částí napravo od $1$, která tvoří graf absolutní hodnoty. Průsečík vyšel na čárkované části, která do grafu nepatří:
\begin{center}
  \begin{tikzpicture}[scale=0.42]
    \draw[->] (-6.8,0) -- (4.5,0) node[right] {\small $x$};
    \draw[->] (0,-7.6) -- (0,6) node[above] {\small $y$};
    \foreach \x in {-6,...,3} \draw (\x,0.15) -- (\x,-0.15);
    \draw[dashed, gray, thick] (1,0) -- (-5.8,-6.8);
    \draw[very thick, colorofwisdom] (-4.2,5.2) -- (1,0) -- (4,3) node[right, font=\small] {$y = |x - 1|$};
    \draw[thick] (-5.6,-7.2) -- (0.6,5.2) node[right, font=\small] {$y = 2x + 4$};
    \filldraw (-1,2) circle (5pt);
    \draw[gray] (-1,2) -- (-2.6,1.2) node[left, font=\small, black] {$(-1, 2)$};
    \filldraw[fill=white, draw=black, thick] (-5,-6) circle (5pt) node[left=3pt, font=\small] {$(-5, -6)$};
    \node[gray, font=\scriptsize, right] at (-2.6,-3.9) {$y = x - 1$};
    \node[font=\small, below] at (1,-0.1) {$1$};
  \end{tikzpicture}
\end{center}

\begin{example}
  Nerovnice se řeší úplně stejně. Vyřešme $|x - 3| - |x| < 1$. Nulové body $0$ a $3$.
  \begin{itemize}
    \item Pro $x < 0$: $-(x - 3) - (-x) = 3 < 1$. Neplatí, $K_1 = \emptyset$.
    \item Pro $0\leq x < 3$: $-(x - 3) - x = 3 - 2x < 1$, tedy $x > 1$. Průnik s intervalem: $K_2 = (1, 3)$.
    \item Pro $x\geq 3$: $(x - 3) - x = -3 < 1$. Platí vždy, $K_3 = [3, \infty)$.
  \end{itemize}
  $K = K_1\cup K_2\cup K_3 = (1, \infty)$. U nerovnic neděláme `kontrolu, zda řešení leží v intervalu', ale rovnou děláme průnik s intervalem.
\end{example}

\begin{note}
  Někdy se vyplatí oběma stranám umocnit. Pokud jsou obě strany nezáporné, je to ekvivalentní úprava. Například
  \begin{equation*}
    |x - 1| > |x + 3|\quad\Leftrightarrow\quad{(x - 1)}^2 > {(x + 3)}^2\quad\Leftrightarrow\quad -2x + 1 > 6x + 9\quad\Leftrightarrow\quad x < -1.
  \end{equation*}
  Geometricky hledáme body, které jsou dál od $1$ než od $-3$. To jsou všechny body vlevo od středu, tedy od $-1$.
  \begin{center}
    \begin{tikzpicture}[scale=0.8]
      \draw[->, thick] (-5.5,0) -- (3.5,0) node[right] {$x$};
      \foreach \x in {-5,...,3} \draw (\x,0.08) -- (\x,-0.08) node[below, font=\scriptsize] {$\x$};
      \draw[line width=3pt, colorofwisdom] (-5.3,0) -- (-1,0);
      \filldraw[fill=white, draw=colorofwisdom, thick] (-1,0) circle (3pt);
      \filldraw (-3,0) circle (2.5pt);
      \filldraw (1,0) circle (2.5pt);
      \draw[dashed, gray] (-1,-0.7) -- (-1,1.3) node[above, font=\scriptsize] {střed};
      \draw[decorate, decoration={brace, amplitude=4pt}] (-3,0.25) -- (-1,0.25) node[midway, above=3pt, font=\scriptsize] {$2$};
      \draw[decorate, decoration={brace, amplitude=4pt}] (-1,0.25) -- (1,0.25) node[midway, above=3pt, font=\scriptsize] {$2$};
      \node[colorofwisdom, font=\small] at (-3.8,-0.9) {blíž k $-3$ než k $1$};
    \end{tikzpicture}
  \end{center}
\end{note}

\begin{problem}
  Vyřešte v $\mathbb{R}$:
  \begin{enumerate}[label=(\alph*)]
    \item $|3x + 2| = 7$,
    \item $|x - 4| < 2x + 1$,
    \item $|x| + |x - 4| = 4$,
    \item $|x + 2| - |x - 1| = x$,
    \item $|x^2 - 4| = 3x$.
  \end{enumerate}
\end{problem}

\begin{solution}
  \begin{enumerate}[label=(\alph*)]
    \item $3x + 2 = \pm 7$, tedy $x = \frac{5}{3}$ nebo $x = -3$.
    \item Pro $x < 4$: $4 - x < 2x + 1$, tedy $x > 1$, průnik $(1, 4)$. Pro $x\geq 4$: $x - 4 < 2x + 1$, tedy $x > -5$, průnik $[4, \infty)$. $K = (1, \infty)$.
    \item Pro $x < 0$: $-2x + 4 = 4$, tedy $x = 0$, neleží v intervalu. Pro $0\leq x < 4$: $4 = 4$, platí vždy. Pro $x\geq 4$: $2x - 4 = 4$, tedy $x = 4$, leží v intervalu. $K = [0, 4]$. Geometricky: součet vzdáleností od $0$ a $4$ je roven $4$ přesně pro body mezi nimi.
    \item Nulové body $-2$ a $1$. Pro $x < -2$: $-x - 2 + x - 1 = -3 = x$, tedy $x = -3$, leží. Pro $-2\leq x < 1$: $x + 2 + x - 1 = 2x + 1 = x$, tedy $x = -1$, leží. Pro $x\geq 1$: $x + 2 - x + 1 = 3 = x$, leží. $K = \{-3, -1, 3\}$.
    \item Pravá strana musí být nezáporná, tedy $x\geq 0$. Pro $0\leq x < 2$: $4 - x^2 = 3x$, $x^2 + 3x - 4 = 0$, $x\in\{1, -4\}$, v intervalu je jen $1$. Pro $x\geq 2$: $x^2 - 4 = 3x$, $x^2 - 3x - 4 = 0$, $x\in\{4, -1\}$, v intervalu je jen $4$. $K = \{1, 4\}$.
  \end{enumerate}
\end{solution}

\subsection{Absolutní hodnota s parametrem}

\begin{example}
  Určeme počet řešení rovnice $|x - 2| = a$ v závislosti na parametru $a\in\mathbb{R}$. Podle geometrického významu: pro $a < 0$ žádné, pro $a = 0$ jedno ($x = 2$), pro $a > 0$ dvě ($x = 2\pm a$).

  Zajímavější je rovnice $|x - 1| + |x + 1| = a$. Z příkladu~\ref{ex:abs-dve} už tušíme, že součet vzdáleností od $-1$ a $1$ je vždy alespoň $2$ (vzdálenost těch dvou bodů), a to přesně pro body mezi nimi. Tedy pro $a < 2$ žádné řešení, pro $a = 2$ celý interval $[-1, 1]$ a pro $a > 2$ dvě řešení $x = \pm\frac{a}{2}$. Kdo nevěří, ať si to rozepíše metodou nulových bodů.
\end{example}

Opět pomůže obrázek. Graf levé strany je `vana' s vodorovným dnem ve výšce $2$ a pravá strana je vodorovná přímka ve výšce $a$. Pokud je přímka pod dnem, nic neprotne. Pokud leží přesně na dně, splyne s ním na celém intervalu $[-1, 1]$. A pokud je výš, protne stěny vany ve dvou bodech:
\begin{center}
  \begin{tikzpicture}[xscale=0.9, yscale=0.55]
    \draw[->] (-3.2,0) -- (3.4,0) node[right] {\small $x$};
    \draw[->] (0,-0.5) -- (0,5.8) node[above] {\small $y$};
    \foreach \x in {-2,-1,1,2} \draw (\x,0.12) -- (\x,-0.12) node[below, font=\scriptsize] {$\x$};
    \draw[very thick] (-2.5,5) -- (-1,2) -- (1,2) -- (2.5,5);
    \draw[thick, colorofwisdom] (-3,3) -- (3,3) node[right, font=\small] {$a = 3$: dvě řešení};
    \filldraw[colorofwisdom] (-1.5,3) circle (3pt);
    \filldraw[colorofwisdom] (1.5,3) circle (3pt);
    \draw[line width=3pt, colorofwisdom!50] (-1,2) -- (1,2);
    \draw[thick, colorofwisdom] (-3,2) -- (3,2) node[right, font=\small] {$a = 2$: celý interval $[-1, 1]$};
    \draw[thick, gray] (-3,1) -- (3,1) node[right, font=\small, black] {$a = 1$: žádné řešení};
  \end{tikzpicture}
\end{center}

\section{Iracionální rovnice}

\begin{definitionbox}
  \textbf{Iracionální rovnice} je rovnice, ve které se neznámá vyskytuje pod odmocninou.

  \vspace{0.2cm}

  \textbf{Postup řešení}:
  \begin{enumerate}
    \item Určíme definiční obor (výrazy pod sudými odmocninami $\geq 0$).
    \item Odmocninu \textbf{osamostatníme} na jedné straně rovnice.
    \item Obě strany \textbf{umocníme}. To je \textbf{důsledková} úprava!
    \item Pokud v rovnici zbyly odmocniny, opakujeme kroky 2 a 3.
    \item Vyřešíme vzniklou rovnici.
    \item Provedeme \textbf{zkoušku}. Ta je zde \textbf{nutnou součástí řešení}, ne jen kontrolou.
  \end{enumerate}
\end{definitionbox}

Proč je zkouška nutná? V kapitole~\ref{kap:rovnice} jsme si řekli, že umocnění není ekvivalentní úprava. Z $a = b$ sice plyne $a^2 = b^2$, ale z $a^2 = b^2$ plyne jen $a = \pm b$. Umocněním tedy do rovnice `propašujeme' i řešení rovnice $a = -b$, která nás vůbec nezajímala. Takovým řešením se říká \textbf{falešné kořeny} a zkouška je jediný spolehlivý způsob, jak je odhalit.

\begin{example}
  Vyřešme $\sqrt{x + 3} = x + 1$. Definiční obor: $x\geq -3$. Umocníme:
  \begin{equation*}
    x + 3 = x^2 + 2x + 1\quad\Leftrightarrow\quad x^2 + x - 2 = 0\quad\Leftrightarrow\quad x\in\{1, -2\}.
  \end{equation*}
  Zkouška:
  \begin{itemize}
    \item $x = 1$: $L = \sqrt{4} = 2$, $P = 2$. Vyhovuje.
    \item $x = -2$: $L = \sqrt{1} = 1$, $P = -1$. \textbf{Nevyhovuje!}
  \end{itemize}
  $K = \{1\}$. Kořen $-2$ je falešný. Přišel z rovnice $\sqrt{x + 3} = -(x + 1)$, kterou jsme umocněním omylem vyřešili také.
\end{example}

Na obrázku je to krásně vidět. Rovnici $x + 3 = {(x + 1)}^2$, kterou jsme po umocnění řešili, splňují body přímky $y = x + 1$, pro které platí $y^2 = x + 3$. To je ale `ležatá' parabola, která má dvě poloviny: horní $y = \sqrt{x + 3}$ (graf naší levé strany) a dolní $y = -\sqrt{x + 3}$. Umocněním jsme k odmocnině nepozorovaně přibrali i její zrcadlový obraz. Falešný kořen $-2$ je průsečík s touto dolní, čárkovanou polovinou:
\begin{center}
  \begin{tikzpicture}[scale=0.75]
    \draw[->] (-3.8,0) -- (3.6,0) node[right] {\small $x$};
    \draw[->] (0,-2.9) -- (0,3.4) node[above] {\small $y$};
    \draw[very thick, colorofwisdom, domain=-3:3, samples=120] plot (\x, {sqrt(\x+3)}) node[right, font=\small] {$y = \sqrt{x + 3}$};
    \draw[dashed, thick, gray, domain=-3:3, samples=120] plot (\x, {-sqrt(\x+3)}) node[right, font=\small] {$y = -\sqrt{x + 3}$};
    \draw[thick] (-3.6,-2.6) -- (2.3,3.3) node[left, font=\small] {$y = x + 1$};
    \filldraw (1,2) circle (2.5pt) node[below right, font=\small] {$(1, 2)$};
    \draw[gray] (-2,-1) -- (-1.2,-2.3) node[below, font=\small, black] {$(-2, -1)$};
    \filldraw[fill=white, draw=black, thick] (-2,-1) circle (2.5pt);
  \end{tikzpicture}
\end{center}

\begin{note}
  Místo zkoušky můžeme na začátku přidat podmínku, že strana, na které není odmocnina, musí být nezáporná (protože odmocnina je nezáporná). V předchozím příkladu $x + 1\geq 0$, tedy $x\geq -1$, a kořen $-2$ by vypadl rovnou. Za těchto podmínek už je umocnění ekvivalentní úprava. U maturity ale doporučujeme dělat zkoušku tak jako tak. Podmínka se snadno pokazí, zkouška ne.
\end{note}

\begin{example}
  Vyřešme $\sqrt{2x + 3} - \sqrt{x + 1} = 1$. Definiční obor: $x\geq -1$. Osamostatníme jednu odmocninu a umocníme:
  \begin{align*}
    \sqrt{2x + 3} &= 1 + \sqrt{x + 1},\\
    2x + 3 &= 1 + 2\sqrt{x + 1} + x + 1,\\
    x + 1 &= 2\sqrt{x + 1}.
  \end{align*}
  Znovu umocníme: $x^2 + 2x + 1 = 4x + 4$, tedy $x^2 - 2x - 3 = 0$ a $x\in\{3, -1\}$. Zkouška: $x = 3$: $\sqrt{9} - \sqrt{4} = 1$, vyhovuje. $x = -1$: $\sqrt{1} - \sqrt{0} = 1$, vyhovuje. $K = \{-1, 3\}$.

  Tentokrát falešné kořeny nevznikly. Ze zkoušky jsme se to ale dozvěděli až na konci, takže ji rozhodně nebylo zbytečné dělat.
\end{example}

\begin{note}
  Ve třetím řádku předchozího příkladu by vás mohlo lákat vydělit rovnici výrazem $\sqrt{x + 1}$. Pak byste přišli o kořen $x = -1$, pro který je $\sqrt{x + 1} = 0$. Vytýkejte, nedělte. Pořád dokola.
\end{note}

\begin{example}
  Někdy stačí podívat se na definiční obor. Rovnice $\sqrt{x - 5} + \sqrt{2 - x} = 1$ vyžaduje $x\geq 5$ a zároveň $x\leq 2$. To nesplňuje žádné číslo, takže $K = \emptyset$ a nemusíme vůbec nic počítat. Kdo by začal umocňovat, dostal by po chvíli nějaká čísla, zkouškou by je pak pracně vyloučil a měl by pocit dobře vykonané práce. Zbytečně.
\end{example}

\begin{example}
  Substituce funguje i tady. Rovnice $x - 5\sqrt{x} + 6 = 0$ se substitucí $t = \sqrt{x}\geq 0$ změní na $t^2 - 5t + 6 = 0$, tedy $t\in\{2, 3\}$ a $x\in\{4, 9\}$. Zkouška: $4 - 10 + 6 = 0$ a $9 - 15 + 6 = 0$.
\end{example}

\begin{problem}
  Vyřešte v $\mathbb{R}$:
  \begin{enumerate}[label=(\alph*)]
    \item $\sqrt{x + 7} = x - 5$,
    \item $\sqrt{3x - 2} = \sqrt{x + 4}$,
    \item $\sqrt{x + 5} + \sqrt{x} = 5$,
    \item $x + \sqrt{x} = 12$,
    \item $\sqrt{x^2 + 9} = x - 3$.
  \end{enumerate}
\end{problem}

\begin{solution}
  \begin{enumerate}[label=(\alph*)]
    \item $x + 7 = x^2 - 10x + 25$, $x^2 - 11x + 18 = 0$, $x\in\{2, 9\}$. Zkouška: $x = 9$: $4 = 4$. $x = 2$: $3\neq -3$. $K = \{9\}$.
    \item $3x - 2 = x + 4$, $x = 3$. Zkouška: $\sqrt{7} = \sqrt{7}$. $K = \{3\}$.
    \item $\sqrt{x + 5} = 5 - \sqrt{x}$, umocníme: $x + 5 = 25 - 10\sqrt{x} + x$, tedy $10\sqrt{x} = 20$, $\sqrt{x} = 2$, $x = 4$. Zkouška: $3 + 2 = 5$. $K = \{4\}$.
    \item Substituce $t = \sqrt{x}$: $t^2 + t - 12 = 0$, $t\in\{3, -4\}$. Záporné $t$ zahodíme, $x = 9$. $K = \{9\}$.
    \item $x^2 + 9 = x^2 - 6x + 9$, tedy $x = 0$. Zkouška: $L = 3$, $P = -3$. $K = \emptyset$. Podmínka $x - 3\geq 0$ by to řekla rovnou.
  \end{enumerate}
\end{solution}

\section{Iracionální nerovnice}

U nerovnic je situace složitější, protože zkouška nefunguje. Nemůžeme dosadit celý interval. Musíme tedy umocňovat jen tehdy, když víme, že jsou obě strany nezáporné, a ostatní případy rozebrat zvlášť.

\begin{example}
  Vyřešme $\sqrt{x - 1} < 3$. Definiční obor $x\geq 1$. Obě strany jsou nezáporné, takže umocnění je ekvivalentní: $x - 1 < 9$, tedy $x < 10$. S definičním oborem $K = [1, 10)$.
\end{example}

\begin{example}
  Vyřešme $\sqrt{x + 2} > x$. Definiční obor $x\geq -2$. Pravá strana může mít libovolné znaménko, a tak rozebereme dva případy.
  \begin{itemize}
    \item $x\in[-2, 0)$: pravá strana je záporná a levá nezáporná, takže nerovnice platí vždy. $K_1 = [-2, 0)$.
    \item $x\geq 0$: obě strany jsou nezáporné a můžeme umocnit: $x + 2 > x^2$, tedy $x^2 - x - 2 < 0$, neboli $(x - 2)(x + 1) < 0$, a tedy $x\in(-1, 2)$. Průnik s $[0, \infty)$ je $K_2 = [0, 2)$.
  \end{itemize}
  $K = K_1\cup K_2 = [-2, 2)$.
\end{example}

\begin{center}
  \begin{tikzpicture}[scale=0.9]
    \draw[->] (-2.8,0) -- (3.5,0) node[right] {\small $x$};
    \draw[->] (0,-2.5) -- (0,3) node[above] {\small $y$};
    \draw[thick, colorofwisdom, domain=-2:3.2, samples=80] plot (\x, {sqrt(\x+2)}) node[right] {\small $y = \sqrt{x + 2}$};
    \draw[thick, domain=-2.4:2.9] plot (\x, \x) node[above right] {\small $y = x$};
    \filldraw (2,2) circle (2pt);
    \filldraw[colorofwisdom] (-2,0) circle (2pt);
    \draw[line width=3pt, colorofwisdom!50] (-2,-0.05) -- (2,-0.05);
    \node[below] at (2,0) {\small $2$};
    \node[below] at (-2,-0.1) {\small $-2$};
  \end{tikzpicture}
\end{center}

Na obrázku je vidět, že graf odmocniny leží nad přímkou $y = x$ přesně na intervalu $[-2, 2)$. Obrázek je u iracionálních nerovnic velice dobrá kontrola.

\begin{problem}
  Vyřešte v $\mathbb{R}$:
  \begin{enumerate}[label=(\alph*)]
    \item $\sqrt{2x + 1}\geq 3$,
    \item $\sqrt{x + 6} < x$,
    \item $\sqrt{4 - x^2} > -1$.
  \end{enumerate}
\end{problem}

\begin{solution}
  \begin{enumerate}[label=(\alph*)]
    \item $x\geq -\frac{1}{2}$, umocníme: $2x + 1\geq 9$, $x\geq 4$. $K = [4, \infty)$.
    \item Definiční obor $x\geq -6$. Pro $x < 0$ je pravá strana záporná a nerovnice neplatí. Pro $x\geq 0$ umocníme: $x + 6 < x^2$, $(x - 3)(x + 2) > 0$, tedy $x > 3$ (průnik s $x\geq 0$). $K = (3, \infty)$.
    \item Odmocnina je vždy nezáporná, takže je větší než $-1$ všude, kde je definovaná. $K = [-2, 2]$. Opět žádné počítání.
  \end{enumerate}
\end{solution}

\section{Odbočka --- Taxíková geometrie}

Absolutní hodnota měří vzdálenost na přímce. Jak měřit vzdálenost v rovině? Obvykle podle Pythagorovy věty: vzdálenost bodů $A = (a_1, a_2)$ a $B = (b_1, b_2)$ je
\begin{equation*}
  d(A, B) = \sqrt{{(a_1 - b_1)}^2 + {(a_2 - b_2)}^2}.
\end{equation*}
Představte si ale, že jste v městě s pravoúhlou sítí ulic, jako je Manhattan. Diagonálně přes domy projít nemůžete, takže skutečná vzdálenost, kterou ujdete, je
\begin{equation*}
  d_{\text{taxi}}(A, B) = |a_1 - b_1| + |a_2 - b_2|.
\end{equation*}
Tomu se říká \textbf{taxíková} (také manhattanská) \textbf{vzdálenost}.\footnote{V Praze to nefunguje, protože ulice nejsou pravoúhlé a taxikář vás stejně poveze oklikou.} Pro body $A = (0, 0)$ a $B = (4, 3)$ je taxíková vzdálenost $4 + 3 = 7$, zatímco vzdušnou čarou je to jen $5$. Zajímavé je, že \textit{každá} cesta po ulicích, která jde jen doprava a nahoru, je stejně dlouhá:
\begin{center}
  \begin{tikzpicture}[scale=0.8]
    \draw[gray!50] (-0.3,-0.3) grid (4.3,3.3);
    \draw[dashed, thick] (0,0) -- (4,3) node[pos=0.45, above left, font=\small] {$5$};
    \draw[very thick, colorofwisdom] (0,0) -- (4,0) -- (4,3);
    \draw[very thick, black!70] (0,0) -- (1,0) -- (1,1) -- (2,1) -- (2,2) -- (3,2) -- (3,3) -- (4,3);
    \filldraw (0,0) circle (2.5pt) node[below left] {$A$};
    \filldraw (4,3) circle (2.5pt) node[above right] {$B$};
    \node[right, font=\small, align=left] at (4.8,1.5) {obě plné cesty\\ mají délku $7$};
  \end{tikzpicture}
\end{center}
Taxíková vzdálenost splňuje všechny rozumné vlastnosti vzdálenosti, včetně trojúhelníkové nerovnosti (dokažte si to pomocí té naší). Geometrie se ale chová velice zvláštně. Co je třeba `kružnice' se středem v počátku a poloměrem $1$? Je to množina bodů, pro které $|x| + |y| = 1$. Metodou nulových bodů (pro každý kvadrant zvlášť) zjistíte, že to je čtverec postavený na špičku:
\begin{center}
  \begin{tikzpicture}[scale=1.3]
    \draw[->] (-1.6,0) -- (1.6,0) node[right] {\small $x$};
    \draw[->] (0,-1.4) -- (0,1.5) node[above] {\small $y$};
    \draw[thick, colorofwisdom] (1,0) -- (0,1) -- (-1,0) -- (0,-1) -- cycle;
    \draw[dashed] (0,0) circle (1);
    \node[below right] at (1,0) {\small $1$};
    \node[above right] at (0,1) {\small $1$};
    \node[colorofwisdom, anchor=south west] at (0.6,0.6) {\small $|x| + |y| = 1$};
  \end{tikzpicture}
\end{center}
Čárkovaně je pro srovnání obyčejná kružnice. Taxíková geometrie není jen hříčka. Používá se v robotice, při zpracování obrazu i ve strojovém učení, kde se mezi `vzdálenostmi' vybírá podle toho, která se pro danou úlohu hodí víc.

\section{Shrnutí a otázky k ústní zkoušce}

\begin{problem}
  Odpovězte vlastními slovy:
  \begin{enumerate}
    \item Jak je definována absolutní hodnota a jaký má geometrický význam?
    \item Jaké má absolutní hodnota vlastnosti? Dokažte trojúhelníkovou nerovnost.
    \item Jak se řeší rovnice a nerovnice tvaru $|x - a| = b$, $|x - a| < b$ a $|x - a| > b$?
    \item Popište metodu nulových bodů. Proč musíme kontrolovat, zda řešení leží v daném intervalu?
    \item Co je iracionální rovnice a jak se řeší? Proč je zkouška nutnou součástí řešení?
    \item Co jsou falešné kořeny a odkud se berou?
    \item Jak se řeší iracionální nerovnice a proč u nich nemůžeme spoléhat na zkoušku?
  \end{enumerate}
\end{problem}
